\section{Introduction}
\label{sec:intro}

The Tsetlin Machine (TM) is a propositional, clause-based learner whose predictions are
assembled from human-readable conjunctions of literals, offering interpretability that
most numeric regressors lack. Its regression variant, the Regression Tsetlin Machine
(RegressionTM), predicts a continuous target from a weighted vote of such clauses. This
transparency comes at a cost familiar to interpretable models: accurately fitting a
regression surface generally requires densely sampled training data, and performance
degrades where the input space is only sparsely covered.

This motivates a data-efficiency question. When training data is scarce or unevenly
distributed, can a better-informed \emph{teacher} model---one trained on more or
better-covered data---guide a data-starved \emph{student} to a better fit? Knowledge
distillation is the standard framing for this question, but distillation for the TM is
under-explored, and the usual output-matching (soft-label) recipe does not obviously map
onto the TM's discrete, clause-level learning dynamics.

We take a different route. Rather than matching the teacher's output, we let the teacher
intervene \emph{inside the student's training loop}, at the exact point where the TM
decides how to reinforce or erode each clause. This makes distillation a modification of
the clause-feedback signal itself, keeping the whole mechanism inside the interpretable
substrate of the model. Our central questions are not only \emph{whether} such guidance
helps, but \emph{when} (which datasets), \emph{where} (which regions of the input space),
and \emph{why}.

We answer these through a sequence of controlled experiments under a masked-distillation
protocol, and arrive at a single unifying account. Our contributions are:

\begin{enumerate}
  \item \textbf{A family of teacher-guided clause-feedback mechanisms} for RegressionTM
    ---a same-side error blend, an adaptive per-sample trigger, opposite-side
    feedback-type overrides, and teacher-bounded update magnitudes---all expressed as
    small, reversible modifications to the standard fit loop.
  \item \textbf{A ``sharpener, not fixer'' characterization:} every mechanism scales
    whatever signal the teacher carries, amplifying a good teacher and aggravating a bad
    one; none manufactures signal a poor teacher lacks.
  \item \textbf{The finding that teacher accuracy does not predict distillation benefit},
    with a clean counterexample (the teacher with the largest standalone accuracy edge
    produces the worst distillation outcome).
  \item \textbf{A demonstration that a 70\%-data guided student can match, and sometimes
    beat, a full-data model}---the latter on datasets where additional data is actively
    harmful.
  \item \textbf{The identification of data spacing/density as the governing variable},
    and a unified mechanism: distillation is a predominantly \emph{global} reshaping of
    the learned function, with a localized hole-filling component that appears only for a
    sufficiently accurate teacher.
\end{enumerate}
